% APPENDIX-ONLY MATERIAL for Table 5.
% Use in the appendix with: % APPENDIX-ONLY MATERIAL for Table 5.
% Use in the appendix with: % APPENDIX-ONLY MATERIAL for Table 5.
% Use in the appendix with: % APPENDIX-ONLY MATERIAL for Table 5.
% Use in the appendix with: \input{table5_appendix_protocol.tex}

\subsection{Table 5 evaluation protocol}
\label{app:table5_protocol}

Table~\ref{tab:table5_simulation_readiness} evaluates the same fixed set of
eight articulated-asset datasets used throughout the study. The reported
denominator is the dataset cohort shown in each panel. In particular, the
free-standing center-of-mass panel uses an eligibility-filtered denominator,
so its $N$ can differ from the nominal 200 assets. In Table~5a-II,
$N_{\mathrm{phys}}$ denotes the number of eligible assets with valid
dataset-provided physics.

\paragraph{Genesis readiness (Table 5a-I).}
Each asset is reset three times to its canonical URDF neutral state. Bounded
coordinates are clamped to their valid ranges during reset. The asset is then
simulated in Genesis for 10 seconds at 240 Hz with gravity, contacts,
self-collision, manifest-bound physics, a fixed root, and zero applied joint
force. Import and DoF mapping require finite simulator states and consistent
canonical mappings. Genesis simulation validity additionally requires finite
states and observed poses at all 2,400 steps in every reset repetition.
Trajectory coverage is the fraction of candidate articulated units with a
complete observed trajectory.

Complete non-placeholder inertials require every required non-root dynamic link
to have a positive finite mass, a finite center of mass, and a valid
positive-definite inertia satisfying the rigid-body triangle inequality. An
asset is not counted as complete when all required links use the exact unit
placeholder inertia convention.

\paragraph{Free-root CoM stability (Table 5a-II).}
For this panel, the root is free and all movable joints are hard-locked at
canonical $q=0$. The initial pose is placed with the collision bounding-box
bottom aligned to an infinite plane, with a 2 mm initial clearance. Gravity is
$[0,0,-9.81]\,\mathrm{m/s^2}$, the timestep is $1/240$ s, and the finite
horizon is 10 seconds. During the first 0.5 seconds, ground-contact positions
within the contact tolerance are deduplicated and used to construct a convex
support polygon in the ground plane.

Physics Coverage is $N_{\mathrm{phys}}/N$, where valid physics requires finite
mass and center-of-mass parameters. CoM Support Coverage is the fraction of
$N_{\mathrm{phys}}$ assets with a non-degenerate contact support polygon. CoM
Static Stability is also computed with $N_{\mathrm{phys}}$ as its denominator;
it requires a non-negative margin (within the prescribed tolerance), an
available support polygon, and a passing full 10-second rollout. If
$N_{\mathrm{phys}}=0$, the conditional CoM metrics are reported as ``N/A''
rather than as physical-instability failures. The signed CoM support margin is
the distance from the horizontal center-of-mass projection to the
support-polygon boundary, positive inside the polygon and negative outside;
the normalized form divides by the support-polygon diameter. The rollout gate
limits transient tilt to 60 degrees, final-window tilt to 15 degrees, and root
drop to one object bounding-box height.

The CoM panel is a diagnostic physical-stability evaluation. To make geometry
usable across heterogeneous released assets, the runtime applies the declared
robust collision policy, including visual-geometry fallback and repair of
degenerate collision meshes when needed, and recomputes inertias for simulator
execution. Dataset-attributable conditional metrics exclude rows lacking valid
source physics, while Physics Coverage reports that missingness explicitly.
Genesis-native fallback mass and CoM are retained only in the supplementary
diagnostic view. This distinction separates physical-parameter coverage from
simulator executability.

\paragraph{Cross-simulator evaluation (Table 5b).}
The same canonical asset selection and joint mapping are evaluated in Genesis,
PyBullet, and MuJoCo. For each asset or joint, an All-3 result is counted only
when the corresponding criterion passes in all three simulators. All-3 Import,
All-3 DoF Mapping, and All-3 Stable are therefore exact intersections, not
averages of three simulator-level percentages. The stable criterion is the
finite-horizon numerical-stability gate from Table 5a-I and should not be
interpreted as physical settling.

\paragraph{Failure and reporting policy.}
Worker errors, incomplete horizons, missing mappings, unavailable contact
support, and missing attributable physics are not silently converted to
passes. Detailed per-asset records, protocol hashes, collision-repair receipts,
and failure categories are retained with the released experiment artifacts.


\subsection{Table 5 evaluation protocol}
\label{app:table5_protocol}

Table~\ref{tab:table5_simulation_readiness} evaluates the same fixed set of
eight articulated-asset datasets used throughout the study. The reported
denominator is the dataset cohort shown in each panel. In particular, the
free-standing center-of-mass panel uses an eligibility-filtered denominator,
so its $N$ can differ from the nominal 200 assets. In Table~5a-II,
$N_{\mathrm{phys}}$ denotes the number of eligible assets with valid
dataset-provided physics.

\paragraph{Genesis readiness (Table 5a-I).}
Each asset is reset three times to its canonical URDF neutral state. Bounded
coordinates are clamped to their valid ranges during reset. The asset is then
simulated in Genesis for 10 seconds at 240 Hz with gravity, contacts,
self-collision, manifest-bound physics, a fixed root, and zero applied joint
force. Import and DoF mapping require finite simulator states and consistent
canonical mappings. Genesis simulation validity additionally requires finite
states and observed poses at all 2,400 steps in every reset repetition.
Trajectory coverage is the fraction of candidate articulated units with a
complete observed trajectory.

Complete non-placeholder inertials require every required non-root dynamic link
to have a positive finite mass, a finite center of mass, and a valid
positive-definite inertia satisfying the rigid-body triangle inequality. An
asset is not counted as complete when all required links use the exact unit
placeholder inertia convention.

\paragraph{Free-root CoM stability (Table 5a-II).}
For this panel, the root is free and all movable joints are hard-locked at
canonical $q=0$. The initial pose is placed with the collision bounding-box
bottom aligned to an infinite plane, with a 2 mm initial clearance. Gravity is
$[0,0,-9.81]\,\mathrm{m/s^2}$, the timestep is $1/240$ s, and the finite
horizon is 10 seconds. During the first 0.5 seconds, ground-contact positions
within the contact tolerance are deduplicated and used to construct a convex
support polygon in the ground plane.

Physics Coverage is $N_{\mathrm{phys}}/N$, where valid physics requires finite
mass and center-of-mass parameters. CoM Support Coverage is the fraction of
$N_{\mathrm{phys}}$ assets with a non-degenerate contact support polygon. CoM
Static Stability is also computed with $N_{\mathrm{phys}}$ as its denominator;
it requires a non-negative margin (within the prescribed tolerance), an
available support polygon, and a passing full 10-second rollout. If
$N_{\mathrm{phys}}=0$, the conditional CoM metrics are reported as ``N/A''
rather than as physical-instability failures. The signed CoM support margin is
the distance from the horizontal center-of-mass projection to the
support-polygon boundary, positive inside the polygon and negative outside;
the normalized form divides by the support-polygon diameter. The rollout gate
limits transient tilt to 60 degrees, final-window tilt to 15 degrees, and root
drop to one object bounding-box height.

The CoM panel is a diagnostic physical-stability evaluation. To make geometry
usable across heterogeneous released assets, the runtime applies the declared
robust collision policy, including visual-geometry fallback and repair of
degenerate collision meshes when needed, and recomputes inertias for simulator
execution. Dataset-attributable conditional metrics exclude rows lacking valid
source physics, while Physics Coverage reports that missingness explicitly.
Genesis-native fallback mass and CoM are retained only in the supplementary
diagnostic view. This distinction separates physical-parameter coverage from
simulator executability.

\paragraph{Cross-simulator evaluation (Table 5b).}
The same canonical asset selection and joint mapping are evaluated in Genesis,
PyBullet, and MuJoCo. For each asset or joint, an All-3 result is counted only
when the corresponding criterion passes in all three simulators. All-3 Import,
All-3 DoF Mapping, and All-3 Stable are therefore exact intersections, not
averages of three simulator-level percentages. The stable criterion is the
finite-horizon numerical-stability gate from Table 5a-I and should not be
interpreted as physical settling.

\paragraph{Failure and reporting policy.}
Worker errors, incomplete horizons, missing mappings, unavailable contact
support, and missing attributable physics are not silently converted to
passes. Detailed per-asset records, protocol hashes, collision-repair receipts,
and failure categories are retained with the released experiment artifacts.


\subsection{Table 5 evaluation protocol}
\label{app:table5_protocol}

Table~\ref{tab:table5_simulation_readiness} evaluates the same fixed set of
eight articulated-asset datasets used throughout the study. The reported
denominator is the dataset cohort shown in each panel. In particular, the
free-standing center-of-mass panel uses an eligibility-filtered denominator,
so its $N$ can differ from the nominal 200 assets. In Table~5a-II,
$N_{\mathrm{phys}}$ denotes the number of eligible assets with valid
dataset-provided physics.

\paragraph{Genesis readiness (Table 5a-I).}
Each asset is reset three times to its canonical URDF neutral state. Bounded
coordinates are clamped to their valid ranges during reset. The asset is then
simulated in Genesis for 10 seconds at 240 Hz with gravity, contacts,
self-collision, manifest-bound physics, a fixed root, and zero applied joint
force. Import and DoF mapping require finite simulator states and consistent
canonical mappings. Genesis simulation validity additionally requires finite
states and observed poses at all 2,400 steps in every reset repetition.
Trajectory coverage is the fraction of candidate articulated units with a
complete observed trajectory.

Complete non-placeholder inertials require every required non-root dynamic link
to have a positive finite mass, a finite center of mass, and a valid
positive-definite inertia satisfying the rigid-body triangle inequality. An
asset is not counted as complete when all required links use the exact unit
placeholder inertia convention.

\paragraph{Free-root CoM stability (Table 5a-II).}
For this panel, the root is free and all movable joints are hard-locked at
canonical $q=0$. The initial pose is placed with the collision bounding-box
bottom aligned to an infinite plane, with a 2 mm initial clearance. Gravity is
$[0,0,-9.81]\,\mathrm{m/s^2}$, the timestep is $1/240$ s, and the finite
horizon is 10 seconds. During the first 0.5 seconds, ground-contact positions
within the contact tolerance are deduplicated and used to construct a convex
support polygon in the ground plane.

Physics Coverage is $N_{\mathrm{phys}}/N$, where valid physics requires finite
mass and center-of-mass parameters. CoM Support Coverage is the fraction of
$N_{\mathrm{phys}}$ assets with a non-degenerate contact support polygon. CoM
Static Stability is also computed with $N_{\mathrm{phys}}$ as its denominator;
it requires a non-negative margin (within the prescribed tolerance), an
available support polygon, and a passing full 10-second rollout. If
$N_{\mathrm{phys}}=0$, the conditional CoM metrics are reported as ``N/A''
rather than as physical-instability failures. The signed CoM support margin is
the distance from the horizontal center-of-mass projection to the
support-polygon boundary, positive inside the polygon and negative outside;
the normalized form divides by the support-polygon diameter. The rollout gate
limits transient tilt to 60 degrees, final-window tilt to 15 degrees, and root
drop to one object bounding-box height.

The CoM panel is a diagnostic physical-stability evaluation. To make geometry
usable across heterogeneous released assets, the runtime applies the declared
robust collision policy, including visual-geometry fallback and repair of
degenerate collision meshes when needed, and recomputes inertias for simulator
execution. Dataset-attributable conditional metrics exclude rows lacking valid
source physics, while Physics Coverage reports that missingness explicitly.
Genesis-native fallback mass and CoM are retained only in the supplementary
diagnostic view. This distinction separates physical-parameter coverage from
simulator executability.

\paragraph{Cross-simulator evaluation (Table 5b).}
The same canonical asset selection and joint mapping are evaluated in Genesis,
PyBullet, and MuJoCo. For each asset or joint, an All-3 result is counted only
when the corresponding criterion passes in all three simulators. All-3 Import,
All-3 DoF Mapping, and All-3 Stable are therefore exact intersections, not
averages of three simulator-level percentages. The stable criterion is the
finite-horizon numerical-stability gate from Table 5a-I and should not be
interpreted as physical settling.

\paragraph{Failure and reporting policy.}
Worker errors, incomplete horizons, missing mappings, unavailable contact
support, and missing attributable physics are not silently converted to
passes. Detailed per-asset records, protocol hashes, collision-repair receipts,
and failure categories are retained with the released experiment artifacts.


\subsection{Table 5 evaluation protocol}
\label{app:table5_protocol}

Table~\ref{tab:table5_simulation_readiness} evaluates the same fixed set of
eight articulated-asset datasets used throughout the study. The reported
denominator is the dataset cohort shown in each panel. In particular, the
free-standing center-of-mass panel uses an eligibility-filtered denominator,
so its $N$ can differ from the nominal 200 assets. In Table~5a-II,
$N_{\mathrm{phys}}$ denotes the number of eligible assets with valid
dataset-provided physics.

\paragraph{Genesis readiness (Table 5a-I).}
Each asset is reset three times to its canonical URDF neutral state. Bounded
coordinates are clamped to their valid ranges during reset. The asset is then
simulated in Genesis for 10 seconds at 240 Hz with gravity, contacts,
self-collision, manifest-bound physics, a fixed root, and zero applied joint
force. Import and DoF mapping require finite simulator states and consistent
canonical mappings. Genesis simulation validity additionally requires finite
states and observed poses at all 2,400 steps in every reset repetition.
Trajectory coverage is the fraction of candidate articulated units with a
complete observed trajectory.

Complete non-placeholder inertials require every required non-root dynamic link
to have a positive finite mass, a finite center of mass, and a valid
positive-definite inertia satisfying the rigid-body triangle inequality. An
asset is not counted as complete when all required links use the exact unit
placeholder inertia convention.

\paragraph{Free-root CoM stability (Table 5a-II).}
For this panel, the root is free and all movable joints are hard-locked at
canonical $q=0$. The initial pose is placed with the collision bounding-box
bottom aligned to an infinite plane, with a 2 mm initial clearance. Gravity is
$[0,0,-9.81]\,\mathrm{m/s^2}$, the timestep is $1/240$ s, and the finite
horizon is 10 seconds. During the first 0.5 seconds, ground-contact positions
within the contact tolerance are deduplicated and used to construct a convex
support polygon in the ground plane.

Physics Coverage is $N_{\mathrm{phys}}/N$, where valid physics requires finite
mass and center-of-mass parameters. CoM Support Coverage is the fraction of
$N_{\mathrm{phys}}$ assets with a non-degenerate contact support polygon. CoM
Static Stability is also computed with $N_{\mathrm{phys}}$ as its denominator;
it requires a non-negative margin (within the prescribed tolerance), an
available support polygon, and a passing full 10-second rollout. If
$N_{\mathrm{phys}}=0$, the conditional CoM metrics are reported as ``N/A''
rather than as physical-instability failures. The signed CoM support margin is
the distance from the horizontal center-of-mass projection to the
support-polygon boundary, positive inside the polygon and negative outside;
the normalized form divides by the support-polygon diameter. The rollout gate
limits transient tilt to 60 degrees, final-window tilt to 15 degrees, and root
drop to one object bounding-box height.

The CoM panel is a diagnostic physical-stability evaluation. To make geometry
usable across heterogeneous released assets, the runtime applies the declared
robust collision policy, including visual-geometry fallback and repair of
degenerate collision meshes when needed, and recomputes inertias for simulator
execution. Dataset-attributable conditional metrics exclude rows lacking valid
source physics, while Physics Coverage reports that missingness explicitly.
Genesis-native fallback mass and CoM are retained only in the supplementary
diagnostic view. This distinction separates physical-parameter coverage from
simulator executability.

\paragraph{Cross-simulator evaluation (Table 5b).}
The same canonical asset selection and joint mapping are evaluated in Genesis,
PyBullet, and MuJoCo. For each asset or joint, an All-3 result is counted only
when the corresponding criterion passes in all three simulators. All-3 Import,
All-3 DoF Mapping, and All-3 Stable are therefore exact intersections, not
averages of three simulator-level percentages. The stable criterion is the
finite-horizon numerical-stability gate from Table 5a-I and should not be
interpreted as physical settling.

\paragraph{Failure and reporting policy.}
Worker errors, incomplete horizons, missing mappings, unavailable contact
support, and missing attributable physics are not silently converted to
passes. Detailed per-asset records, protocol hashes, collision-repair receipts,
and failure categories are retained with the released experiment artifacts.
